校准需要一种不依赖近邻标签的比较对象，这里采用同一训练索引在较大预算下返回的 ID。令 $\mathcal R_\alpha(q)$ 包含 $k$ 个有效、互异的返回 ID，并写成 $\mathcal R_{\rm ref}(q)=\mathcal R_{\alpha_{\rm ref}}(q)$。测试网格和参考预算为
\begin{equation}
 A=(200,100,64,32,16,8,4,2),\qquad \alpha_{\rm ref}=A[0]=200.
 \label{eq:6-grid}
\end{equation}
对同一组 $S$ 个查询 $\mathcal Q_S$，忽略结果顺序和距离，计数小预算未能返回的参考 ID：
\begin{equation}
 \begin{aligned}
 D_S(\alpha)
 &=\sum_{q\in\mathcal Q_S}
   \left(k-|\mathcal R_\alpha(q)\cap\mathcal R_{\rm ref}(q)|\right)\\
 &=\frac12\sum_{q\in\mathcal Q_S}
   |\mathcal R_\alpha(q)\triangle\mathcal R_{\rm ref}(q)|.
 \end{aligned}
 \label{eq:6-agreement}
\end{equation}
接受规则 $D_S(\alpha)\leq\epsilon=1$ 表示整个样本最多缺失一个参考 ID。目标是找到通过该规则的最小网格预算：
\begin{equation}
 \widehat\alpha=\min\{\alpha\in A:D_S(\alpha)\leq\epsilon\}.
 \label{eq:6-selection-objective}
\end{equation}
没有较小预算被接受时，使用参考预算 200。\Cref{sec:6-3} 说明二分何时能精确找到这一最小值。

由于容许量始终只有一个 ID，增大样本会降低允许缺失的参考答案比例。当 $k=10$ 时，
\[
 \begin{array}{c|rrrr}S&32&64&128&500\\\hline
 \epsilon/(Sk)&0.003125&0.0015625&0.00078125&0.0002
 \end{array}
\]
大样本既覆盖更多查询，也施加更严格的相对阈值。\cref{sec:7-5} 中独立的测试查询评估从校准池选出的预算能否保持质量。

\par\medskip\noindent\begin{minipage}{\linewidth}
\captionsetup{hypcap=false}
\centering
% Native schematic of the implemented decision and its offline evaluation.
\begin{tikzpicture}[
 font=\footnotesize\linespread{1}\selectfont,
 box/.style={draw=ThesisNavy!65,rounded corners=2pt,align=center,
 text width=2.8cm,minimum height=1.45cm,inner sep=4pt},
 arr/.style={-{Latex[length=2mm]},draw=ThesisNavy!80,thick}]
\node[box] (sample) at (0,0) {Calibration pool\\Even indices\\Sample $S$ queries};
\node[box,fill=ThesisNavy!7] (ref) at (3.6,0) {Reference search\\$\alpha_{\rm ref}=200$\\Top-$k$ IDs};
\node[box,fill=ThesisNavy!7] (test) at (7.2,0) {Bisection\\$D_S(\alpha)\leq1$\\Shared sample};
\node[box] (choice) at (10.8,0) {Chosen $\alpha$\\Caller applies it\\For batch reuse};
\node[box] (held) at (0,-2.0) {Held-out pool\\Odd query indices\\GT for evaluation};
\node[box,text width=6.4cm] (eval) at (7.2,-2.0)
 {Evaluate the selected arm\\Look up held-out recall and QPS\\from previously timed grid arms};
\draw[arr] (sample) -- (ref);
\draw[arr] (ref) -- (test);
\draw[arr] (test) -- (choice);
\draw[arr] (held) -- (eval);
\draw[arr,dashed] (choice.south) -- (10.8,-2.0) -- (eval.east);
\end{tikzpicture}

\captionof{figure}[工作负载级校准与评估]{工作负载级校准与评估.}
\label{fig:6-calibration}
\end{minipage}\par\medskip
\cref{fig:6-calibration} 中参考搜索和试验搜索使用同一组 $S$ 个查询。GT 在预算选择后、于留出查询上评估质量，所选预算的吞吐来自固定预算测量。\Cref{sec:a-budget-contract} 定义这些计时包含的调用。

\subsection{一致性控制什么}

当保留候选嵌套且最终排序固定时，不一致有直接解释：它计数被主码筛选排除的参考答案。在 \Cref{prop:6-monotonicity} 下，
\begin{equation}
 D_S(\alpha)=\sum_{q\in\mathcal Q_S}
       |\mathcal R_{\rm ref}(q)\setminus C_\alpha(q)|.
 \label{eq:6-reference-coverage}
\end{equation}
对任意含 $k$ 个元素的输出和 GT 集合 $\mathcal T_k$，GT 命中损失不会超过缺失的参考 ID 数：
\[
 |\mathcal R_{\rm ref}\cap\mathcal T_k|-|\mathcal R_\alpha\cap\mathcal T_k|
 \leq|\mathcal R_{\rm ref}\setminus\mathcal R_\alpha|.
\]
在样本上求和再除以 $Sk$，得到
\begin{equation}
 \begin{aligned}
 r_{\mathcal Q_S}(\alpha)
 &\geq r_{\mathcal Q_S}({\rm ref})-\frac{D_S(\alpha)}{Sk},\\
 D_S(\alpha)\leq\epsilon
 &\implies r_{\mathcal Q_S}(\alpha)
       \geq r_{\mathcal Q_S}({\rm ref})-\frac{\epsilon}{Sk}.
 \end{aligned}
 \label{eq:6-relative-recall}
\end{equation}
这是校准样本上相对参考的召回率界。参考本身的误差仍由 GT 评估（\cref{sec:7-recall-attribution}），该界不会自动迁移至未见查询。

\subsection{抽样分辨率与有序检验}

小样本可能漏掉不常发生的返回 ID 变化。下面的模型说明，这些变化能以多高概率仍未被观察到。首先，令 $H_S$ 计数缺失任意参考 ID 的查询数。由于 $H_S\leq D_S$，被接受预算满足 $H_S\leq1$。对预先指定的预算及独立同分布查询，以 $u$ 表示该查询事件概率，容许一次事件的单侧 95\% 二项上限满足 \cite{nist_binomial}
\begin{equation}
 (1-u)^S+Su(1-u)^{S-1}=0.05.
 \label{eq:6-binomial-limit}
\end{equation}
更强的模型将每个返回近邻位置视为概率相同的独立事件。把上限应用到全部结果槽位，得到下表第二行：
\[
 \begin{array}{c|cc}
 &S=128&S=500\\\hline
 \text{i.i.d. query events}&0.03652&0.009452\\
 \text{i.i.d. homogeneous result slots}&0.003701&0.0009484
 \end{array}
\]
槽位模型给出约 $4.74/(Sk)$：$S=500$ 时相当于千分之一左右的参考结果缺失率，$S=128$ 时则约为千分之四。查询独立并不使各结果槽位独立。这些是模型下的参考不一致分辨率，而非重复使用固定池样本的保证。

只有每个预算检验具有有效的统计误差界时，单调性才允许按固定序列检验预算 \cite{learn_then_test}。实现中的一个 ID 接受阈值检查的是一致性，本身不认证未来查询的 GT 召回率。

\paragraph{有限池扩展。}
对实际抽样方式，可以改为从校准池的构成计算接受概率。令 $X_q(\alpha)$ 计数查询 $q$ 缺失的参考 ID，假设它在该执行区间内不随查询进入的子集或位置改变。池大小为 $N_{\rm cal}$，令 $n_j=|\{q:X_q=j\}|$。均匀无放回样本仅在全部查询无缺失，或恰有一个查询缺失一个 ID 时通过 \cite{nist_hypergeometric}：
\begin{equation}
 \Pr\{D_S(\alpha)\leq1\}
 =\frac{\binom{n_0}{S}+n_1\binom{n_0}{S-1}}{\binom{N_{\rm cal}}{S}}.
 \label{eq:6-finite-pool}
\end{equation}
越界二项系数取零。将接受概率转为 GT 损失预测，还需知道哪些预算通过及各预算的留出质量。\Cref{sec:7-budget-trace} 将记录的子集和与这一计算对照。
